% =====================================================================
%  2bat-ud3-11.tex  —  SESSIÓ 11: Explorar amb GeoGebra (asímptotes i discontinuïtats)
% =====================================================================
\horatitol{Sessió 11 — Explorar amb GeoGebra}%
{Tancar el treball tecnològic: localitzar asímptotes i discontinuïtats a la pantalla i validar gràficament les conclusions fetes a mà.}

\subsection{Idea clau}

Tanquem el treball amb GeoGebra ajuntant els dos grans temes de la unitat: les
funcions \textbf{asimptòtiques} (límits) i les definides \textbf{a trossos}
(continuïtat). L'objectiu és \textbf{validar gràficament} el que ja vam calcular a mà:
el càlcul, la gràfica i la interpretació han de \textbf{coincidir}.

\begin{idea}
\textbf{Què validem a GeoGebra}
\begin{itemize}[leftmargin=1.4em, itemsep=2pt]
  \item En una funció asimptòtica: que el \textbf{límit a l'infinit} calculat a mà es
        correspon amb l'\textbf{asímptota horitzontal} que es veu a la pantalla.
  \item En una funció a trossos: que allà on havíem detectat un \textbf{salt}, GeoGebra
        mostra efectivament una \textbf{discontinuïtat} (un trencament a la gràfica).
\end{itemize}
\textbf{Funcions a trossos a GeoGebra:} s'escriuen amb l'ordre \texttt{Si(...)} (o
\texttt{If}), indicant la condició de cada tram. Per exemple:
\texttt{f(x)=Si(x<2, x+1, x+3)}.
\textbf{La regla d'or:} la gràfica i l'anàlisi s'han de validar \textbf{mútuament}; si
no coincideixen, una de les dues té un error.
\end{idea}

\subsection{Exemples}

\begin{exemple}[validar un límit i una asímptota]
Introduïm \texttt{f(x)=600-900/x} i la recta \texttt{y=600}. En ampliar la vista cap a
la dreta, la gràfica s'acosta a $y=600$ sense tocar-la. Això \textbf{valida} el càlcul
a mà: $\lim_{x\to\infty}f(x)=600$, asímptota horitzontal $y=600$.
\end{exemple}

\begin{exemple}[validar una discontinuïtat]
Introduïm la funció a trossos
\[
g(x)=
\begin{cases}
x+1 & \text{si } x\le 2,\\[2pt]
x+3 & \text{si } x>2,
\end{cases}
\qquad\text{(a GeoGebra: \texttt{g(x)=Si(x<=2, x+1, x+3)})}.
\]
A la pantalla es veu un \textbf{trencament} a $x=2$: la gràfica salta de $3$ a $5$.
Això valida el que havíem trobat a mà: discontinuïtat de salt a $x=2$.

\begin{center}
\begin{tikzpicture}
\begin{axis}[udaxis, width=8cm, height=4.4cm, xlabel={\small $x$}, ylabel={\small $y$},
    xmin=-1, xmax=6, ymin=0, ymax=9, xtick={0,2,4}, ytick={0,3,5}]
  \addplot[udblue, domain=-1:2, samples=2]{x+1};
  \addplot[udblue, domain=2:5, samples=2]{x+3};
  \addplot[udnavy, only marks, mark=*, mark size=1.5pt] coordinates {(2,3)};
  \addplot[udnavy, only marks, mark=o, mark size=1.5pt] coordinates {(2,5)};
\end{axis}
\end{tikzpicture}

\figcap{GeoGebra mostra el trencament a $x=2$, validant el salt calculat a mà.}
\end{center}
\end{exemple}

\subsection{Exercicis resolts}

\begin{exresolt}[validar la coherència de les tres mirades]
Has calculat a mà que $f(x)=300-\dfrac{600}{x}$ té asímptota horitzontal $y=300$. Com
ho confirmaries a GeoGebra i què hauries de veure?
\solucio
Introdueixes \texttt{f(x)=300-600/x} i \texttt{y=300}, i amplies la vista cap a la
dreta. Hauries de veure que la gràfica s'acosta a la recta $y=300$ sense traspassar-la.
Les tres mirades coincideixen: \textbf{càlcul} ($\lim=300$), \textbf{gràfica} (s'acosta
a $y=300$) i \textbf{interpretació} (el fenomen s'estabilitza en $300$).
\resposta{La gràfica s'acosta a $y=300$; càlcul, gràfica i interpretació coincideixen.}
\end{exresolt}

\begin{exresolt}[la gràfica corregeix una conclusió]
A mà has dit que $g(x)=\begin{cases}2x & x\le 1\\ x+1 & x>1\end{cases}$ té un salt a
$x=1$. A GeoGebra veus la gràfica \emph{sense} trencament. Qui té raó?
\solucio
Comprovem el càlcul: a $x=1$, esquerre $2\cdot 1=2$; dret $1+1=2$. \textbf{Coincideixen},
així que la funció és \textbf{contínua}: no hi ha salt. GeoGebra té raó; la conclusió
a mà era errònia (probablement per avaluar malament un tram). La gràfica ha ajudat a
corregir-la.
\resposta{No hi ha salt: la funció és contínua. La conclusió a mà era errònia.}
\end{exresolt}

\subsection{Exercicis per fer}

\noindent\textit{Full de pràctica tecnològica. Per a cada funció: anota el que has
trobat a mà, comprova-ho a GeoGebra i digues si coincideix.}

\begin{exercicis}
\begin{exlist}
  \item Introdueix \texttt{f(x)=600-900/x} i \texttt{y=600}. Amplia la vista i comprova
        que la gràfica s'acosta a l'asímptota. Coincideix amb el teu càlcul del límit?
  \item Introdueix \texttt{g(x)=Si(x<=2, x+1, x+3)} i localitza visualment la
        discontinuïtat. A quin valor de $x$ es produeix el salt i de quina mida és?
  \item Introdueix \texttt{h(x)=Si(x<=1, 2x, x+1)} i digues si la gràfica és contínua o
        té salt. Coincideix amb la comprovació a mà (avalua els dos trams a $x=1$)?
  \item Completa el full de validació amb tres funcions (asimptòtiques o a trossos):
        \begin{center}
        \renewcommand{\arraystretch}{1.4}
        \begin{tabular}{p{3.4cm} p{4.4cm} p{2.4cm} p{1.6cm}}
        \toprule
        \textbf{Funció} & \textbf{Conclusió a mà} & \textbf{GeoGebra} & \textbf{Coincideix?} \\
        \midrule
         & & & \\[10pt]
         & & & \\[10pt]
         & & & \\
        \bottomrule
        \end{tabular}
        \end{center}
  \item \textbf{(Validació mútua)} En algun cas, la gràfica t'ha ajudat a corregir una
        conclusió feta a mà? Explica què havia passat.
  \item \textbf{(Síntesi)} Explica per què diem que el càlcul, la gràfica i la
        interpretació «s'han de validar mútuament». Què vol dir que les tres mirades
        coincideixin?
\end{exlist}
\end{exercicis}

% ---------------------------------------------------------------------
\fullsolucions{Sessió 11}
\begin{solucions}
\begin{sollist}
  \item La gràfica s'acosta a $y=600$ sense tocar-la; coincideix amb
        $\lim_{x\to\infty}f(x)=600$.
  \item Salt a $x=2$: el tram esquerre arriba a $3$ i el dret surt de $5$; salt de mida
        $5-3=2$.
  \item A $x=1$: esquerre $2\cdot 1=2$, dret $1+1=2$. Coincideixen: \textbf{contínua},
        sense salt. GeoGebra ho confirma.
  \item Resposta oberta (full de validació personal). Es valora la coherència entre el
        càlcul i la visualització.
  \item Resposta oberta. Es valora descriure bé el cas i l'error corregit.
  \item Resposta oberta. Idea: el càlcul dóna el resultat numèric, la gràfica el mostra
        visualment i la interpretació l'explica en context; si les tres coincideixen, la
        conclusió és sòlida, i si no, hi ha un error en alguna que cal trobar.
\end{sollist}
\end{solucions}
